\documentclass[11pt,a4paper]{article}

\usepackage[T1]{fontenc}
\usepackage[utf8]{inputenc}
\usepackage{newpxtext,newpxmath}
\usepackage{microtype}
\usepackage[a4paper, left=2.35cm, right=2.35cm, top=2.4cm, bottom=2.45cm, headheight=17pt]{geometry}

\usepackage{amsmath,mathtools}
\usepackage{braket}
\usepackage{graphicx}
\usepackage{booktabs}
\usepackage{array}
\usepackage{enumitem}
\usepackage{caption}
\usepackage{xcolor}
\usepackage{titlesec}
\usepackage{fancyhdr}
\usepackage{tocloft}
\usepackage[most]{tcolorbox}
\usepackage[
    colorlinks=true,
    linkcolor=pathBlue,
    citecolor=accentTeal,
    urlcolor=pathBlue,
    pdftitle={Path Integrals in Quantum Mechanics}
]{hyperref}

\allowdisplaybreaks
\setlength{\parindent}{0pt}
\setlength{\parskip}{0.48em}
\linespread{1.06}

% Palette
\definecolor{inkDark}{HTML}{1E293B}
\definecolor{slateSub}{HTML}{334155}
\definecolor{pathBlue}{HTML}{005A9C}
\definecolor{accentTeal}{HTML}{0E6251}
\definecolor{accentAmber}{HTML}{B45309}
\definecolor{tintBlue}{HTML}{F0F7FF}
\definecolor{tintTeal}{HTML}{F0FDF9}
\definecolor{tintAmber}{HTML}{FFFBEB}
\definecolor{ruleGray}{HTML}{CBD5E1}

\color{inkDark}

% Headings & ToC
\titleformat{\section}
    {\Large\bfseries\color{inkDark}}
    {\color{pathBlue}\thesection}{0.7em}{}
    [\vspace{2pt}{\color{ruleGray}\titlerule[0.7pt]}]
\titlespacing*{\section}{0pt}{1.4em}{0.55em}

\titleformat{\subsection}
    {\large\bfseries\color{pathBlue}}
    {\thesubsection}{0.6em}{}
\titlespacing*{\subsection}{0pt}{1.05em}{0.35em}

\setlength{\cftbeforesecskip}{3pt}
\setlength{\cftbeforesubsecskip}{0.5pt}
\renewcommand{\cftaftertoctitle}{\par\vspace{2pt}}

\captionsetup{
    font={small,color=slateSub},
    labelfont={bf,color=inkDark},
    width=0.94\textwidth,
    skip=6pt
}

\pagestyle{fancy}
\fancyhf{}
\fancyhead[L]{\small\color{slateSub}\textsc{Path Integrals in Quantum Mechanics}}
\fancyhead[R]{\small\color{slateSub}\nouppercase{\leftmark}}
\fancyfoot[C]{\small\color{slateSub}\thepage}
\renewcommand{\headrulewidth}{0.5pt}
\renewcommand{\headrule}{\hbox to\headwidth{\color{ruleGray}\leaders\hrule height \headrulewidth\hfill}}

% Boxes
\tcbset{
    notebox/.style={
        enhanced,
        boxrule=0pt,
        leftrule=3pt,
        arc=2pt,
        outer arc=2pt,
        top=6pt,
        bottom=6pt,
        left=9pt,
        right=9pt,
        before skip=0.75em,
        after skip=0.75em,
        fonttitle=\bfseries\normalsize,
        attach title to upper={\par\vspace{2pt}},
        coltext=inkDark
    }
}

\newtcolorbox{definitionbox}[1][]{
    notebox, colback=tintBlue, colframe=pathBlue, coltitle=pathBlue,
    title={Definition\ifstrempty{#1}{}{ --- #1}}
}
\newtcolorbox{theorembox}[1][]{
    notebox, colback=tintTeal, colframe=accentTeal, coltitle=accentTeal,
    title={Theorem\ifstrempty{#1}{}{ --- #1}}
}
\newtcolorbox{propositionbox}[1][]{
    notebox, colback=tintTeal, colframe=accentTeal, coltitle=accentTeal,
    title={Proposition\ifstrempty{#1}{}{ --- #1}}
}
\newtcolorbox{remarkbox}[1][]{
    notebox, colback=tintAmber, colframe=accentAmber, coltitle=accentAmber,
    title={Remark\ifstrempty{#1}{}{ --- #1}}
}

% Macros
\newcommand{\ii}{\mathrm{i}}
\newcommand{\dd}{\mathrm{d}}
\newcommand{\ee}{\mathrm{e}}
\newcommand{\Lag}{\mathcal{L}}
\newcommand{\Dpath}{\mathcal{D}}
\newcommand{\xcl}{x_{\mathrm{cl}}}
\newcommand{\Scl}{S_{\mathrm{cl}}}
\newcommand{\order}{\mathcal{O}}

\begin{document}

\thispagestyle{plain}
\begin{center}
    {\color{pathBlue}\rule{\textwidth}{1.5pt}}\\[0.75em]
    {\LARGE\bfseries From Classical Paths to Quantum Path Integrals}\\[0.35em]
    {\normalsize\color{slateSub} Feynman's Historical Construction, the Operator Proof, and the Free Particle}\\[0.55em]
    {\color{ruleGray}\rule{\textwidth}{0.7pt}}
\end{center}

{\small\setlength{\parskip}{0.12em}\tableofcontents}
\vspace{0.4em}

% ============================================================================
\section[Classical Paths vs.\ Quantum Amplitudes]{Classical Paths vs.\ Quantum Amplitudes}
\label{sec:guiding}
% ============================================================================

In classical mechanics, a one-dimensional particle moving from $(t_i, x_i)$ to $(t_f, x_f)$ follows a curve $x(t):[t_i,t_f]\to\mathbb{R}$. For a Lagrangian $\Lag(x,\dot x,t)=T-V$, the action functional is
\begin{equation}
    S[x(t)] := \int_{t_i}^{t_f} \Lag\bigl(x(t),\dot x(t),t\bigr)\,\dd t.
\end{equation}
The physical trajectory $\xcl(t)$ makes the action stationary, $\delta S[\xcl]=0$, which yields the Euler--Lagrange equation
\begin{equation}
    \frac{\dd}{\dd t}\!\left(\frac{\partial\Lag}{\partial\dot x}\right) = \frac{\partial\Lag}{\partial x}.
    \label{eq:euler-lagrange}
\end{equation}

How does this single path become quantum? A common guess is that quantum uncertainty thickens $\xcl(t)$ into a narrow ``tube.'' That picture works as a semiclassical approximation when $S \gg \hbar$, but it is not the underlying rule:
\begin{enumerate}[leftmargin=1.5em, itemsep=2pt]
    \item \textbf{Sum over paths.} A quantum particle is not confined to a tube. Every path connecting $(t_i,x_i)$ to $(t_f,x_f)$ contributes to the transition amplitude with equal magnitude and phase $\ee^{\ii S[x]/\hbar}$ (Figure~\ref{fig:classical-vs-quantum}). When $S\gg\hbar$, phases vary rapidly away from $\xcl(t)$ and cancel out, leaving an effective constructive region $|S[x]-S[\xcl]|\lesssim\hbar$.
    \item \textbf{Wave-packet evolution.} In Hilbert space, a state $\ket{\psi(t)}$ has position density $|\psi(x,t)|^2$ with finite width. By Ehrenfest's theorem,
    \begin{equation}
        m\frac{\dd^2\langle\hat x\rangle}{\dd t^2} = -\langle V'(\hat x)\rangle.
    \end{equation}
    If $V(x)$ is at most quadratic or varies slowly across the packet, $\langle\hat x\rangle(t)$ follows $\xcl(t)$ while $|\psi(x,t)|^2$ traces a band around it (Section~\ref{sec:free-packet}).
\end{enumerate}

\begin{figure}[htbp]
    \centering
    \includegraphics[width=0.96\textwidth]{assets/fig1_classical_vs_quantum.pdf}
    \caption{\textbf{(a)} Classical mechanics picks the stationary path $\xcl(t)$ ($\delta S=0$). \textbf{(b)} Quantum mechanics sums $\ee^{\ii S[x]/\hbar}$ over all paths; for $S\gg\hbar$, constructive interference concentrates inside $|S[x]-S[\xcl]|\lesssim\hbar$.}
    \label{fig:classical-vs-quantum}
\end{figure}

The ``tube'' picture fails in three standard situations:
\begin{itemize}[leftmargin=1.5em, itemsep=2pt]
    \item \textbf{When $S \sim \hbar$:} the constructive region spreads over the whole setup or splits into separate channels, as in double-slit interference.
    \item \textbf{Destructive interference:} a classical random walk inside a tube adds positive probabilities, whereas quantum paths add complex amplitudes $\ee^{\ii S/\hbar}$ and can cancel to zero.
    \item \textbf{Tunneling ($E < V$):} no real classical trajectory exists inside the barrier, yet the quantum amplitude is nonzero.
\end{itemize}

% ============================================================================
\section[Historical Route: From Lagrangian to Schr\"odinger]{Historical Route: From the Lagrangian to the Schr\"odinger Equation}
\label{sec:historical}
% ============================================================================

Textbooks usually start from $\hat U(T) = \ee^{-\ii\hat H T/\hbar}$ and derive the path integral. Feynman originally worked in the opposite direction: he started from the Lagrangian $\Lag$ and derived the Schr\"odinger equation.

\subsection{Why Feynman Needed a Formulation Without a Hamiltonian}

In 1941--1942, Feynman and Wheeler studied action-at-a-distance electrodynamics, in which charges interact directly at retarded and advanced times $t \mp r/c$ without an independent electromagnetic field. One can write an action $S$ for the particle worldlines, but because different times are coupled directly, there is no single-time Hamiltonian $H$. Standard canonical quantization ($\hat H$, $[\hat x,\hat p]=\ii\hbar$) could not be applied, so Feynman needed a way to quantize directly from $\Lag$ and $S$.

Herbert Jehle pointed Feynman to Dirac's 1933 paper \emph{``The Lagrangian in Quantum Mechanics,''} where Dirac noted from Hamilton--Jacobi theory that the short-time kernel $\braket{x_{t+\varepsilon}|x_t}$ is ``analogous to'' $\exp(\ii\varepsilon\Lag/\hbar)$. Feynman tested whether ``analogous to'' meant strictly proportional on a 1D particle with $\Lag(x,\dot x) = \frac{1}{2}m\dot x^2 - V(x)$.

By linearity, the wave function at time $t+\varepsilon$ ($\varepsilon > 0$ small) is given by an integral kernel:
\begin{equation}
    \psi(x,t+\varepsilon) = \frac{1}{A(\varepsilon)}\int_{-\infty}^{\infty}
    \exp\!\left\{\frac{\ii\varepsilon}{\hbar}\left[
        \frac{m}{2}\!\left(\frac{x-y}{\varepsilon}\right)^2
        - V\!\left(\frac{x+y}{2}\right)
    \right]\right\}\psi(y,t)\,\dd y,
    \label{eq:feynman-ansatz}
\end{equation}
where $A(\varepsilon)$ is a normalization constant to be fixed. No operators or commutators are assumed in~\eqref{eq:feynman-ansatz}.

\subsection{Deriving the Schr\"odinger Equation}

Let $\eta := y - x$ (so $y = x+\eta$ and $\dd y = \dd\eta$). Equation~\eqref{eq:feynman-ansatz} becomes
\begin{equation}
    \psi(x,t+\varepsilon) = \frac{1}{A(\varepsilon)}\int_{-\infty}^{\infty}
    \exp\!\left(\frac{\ii m\eta^2}{2\hbar\varepsilon}\right)
    \exp\!\left[-\frac{\ii\varepsilon}{\hbar}V\!\left(x+\frac{\eta}{2}\right)\right]
    \psi(x+\eta,t)\,\dd\eta.
    \label{eq:eta-integral}
\end{equation}

Define $\sigma_\varepsilon := \sqrt{\hbar\varepsilon/m}$. For $|\eta| \gg \sigma_\varepsilon$, the factor $\exp\bigl(\frac{\ii m\eta^2}{2\hbar\varepsilon}\bigr)$ oscillates rapidly and cancels in the integral (Figure~\ref{fig:fresnel-scaling}). The integral is therefore dominated by $|\eta| \sim \sigma_\varepsilon = \order(\varepsilon^{1/2})$.

\begin{figure}[htbp]
    \centering
    \includegraphics[width=0.94\textwidth]{assets/fig2_fresnel_scaling.pdf}
    \caption{Real part of the kinetic kernel $\exp\bigl(\frac{\ii m\eta^2}{2\hbar\varepsilon}\bigr)$ with $\sigma_\varepsilon=\sqrt{\hbar\varepsilon/m}$. Contributions concentrate in $|\eta|\lesssim\sigma_\varepsilon=\order(\varepsilon^{1/2})$, so $\eta^2$ is of order $\varepsilon$.}
    \label{fig:fresnel-scaling}
\end{figure}

Because $\eta^2 \sim \order(\varepsilon)$, keeping all terms up to $\order(\varepsilon)$ requires expanding $\psi(x+\eta,t)$ to second order in $\eta$:
\begin{align}
    \psi(x+\eta,t) &= \psi(x,t) + \eta\,\partial_x\psi(x,t) + \frac{1}{2}\eta^2\,\partial_x^2\psi(x,t) + \order(\eta^3), \\
    \exp\!\left[-\frac{\ii\varepsilon}{\hbar}V\!\left(x+\frac{\eta}{2}\right)\right]
    &= 1 - \frac{\ii\varepsilon}{\hbar}V(x) - \frac{\ii\varepsilon}{2\hbar}\eta V'(x) + \order(\varepsilon^2).
\end{align}
Here the $-\frac{\ii\varepsilon}{2\hbar}\eta V'(x)$ term is $\order(\varepsilon^{3/2})$ and odd in $\eta$, so evaluating $V$ at $x$, $y$, or $\frac{1}{2}(x+y)$ gives the same result to $\order(\varepsilon)$.

\begin{propositionbox}[Fresnel--Gaussian Moments]
For $\alpha\in\mathbb{C}$ with $\operatorname{Re}\alpha > 0$ (principal root $\operatorname{Re}\sqrt{\alpha}>0$),
\begin{equation}
    \int_{-\infty}^{\infty} \ee^{-\alpha\eta^2}\,\dd\eta = \sqrt{\frac{\pi}{\alpha}}, \qquad
    \int_{-\infty}^{\infty} \eta\,\ee^{-\alpha\eta^2}\,\dd\eta = 0, \qquad
    \int_{-\infty}^{\infty} \eta^2\,\ee^{-\alpha\eta^2}\,\dd\eta = \frac{1}{2\alpha}\sqrt{\frac{\pi}{\alpha}}.
    \label{eq:gaussian-moments}
\end{equation}
With the convergence prescription $\varepsilon \to \varepsilon - \ii 0^+$, setting $\alpha = \frac{m}{2\ii\hbar\varepsilon} + 0^+$ gives
\begin{equation}
    I_0 := \int_{-\infty}^{\infty}\ee^{\frac{\ii m\eta^2}{2\hbar\varepsilon}}\,\dd\eta
    = \sqrt{\frac{2\pi\ii\hbar\varepsilon}{m}}, \qquad
    I_1 = 0, \qquad
    I_2 := \int_{-\infty}^{\infty}\eta^2\,\ee^{\frac{\ii m\eta^2}{2\hbar\varepsilon}}\,\dd\eta
    = \frac{\ii\hbar\varepsilon}{m}\,I_0,
    \label{eq:fresnel-moments}
\end{equation}
where $\sqrt{\ii} = \ee^{\ii\pi/4}$.
\end{propositionbox}

\textit{Proof.}\;
The first integral holds for $\alpha > 0$ via polar coordinates and extends to $\operatorname{Re}\alpha > 0$ by analytic continuation. Differentiating with respect to $-\alpha$ gives the $\eta^2$ moment; the $\eta$ moment vanishes by parity. \hfill$\square$

Substituting the expansions and~\eqref{eq:fresnel-moments} into~\eqref{eq:eta-integral} (the $\eta^3$ term also drops out by parity) gives
\begin{equation}
    \psi(x,t+\varepsilon)
    = \frac{I_0}{A(\varepsilon)}\left[
        \psi(x,t)
        - \frac{\ii\varepsilon}{\hbar}V(x)\psi(x,t)
        + \frac{\ii\hbar\varepsilon}{2m}\,\partial_x^2\psi(x,t)
        + \order(\varepsilon^2)
    \right].
    \label{eq:feynman-expanded}
\end{equation}
Matching orders of $\varepsilon$:
\begin{itemize}[leftmargin=1.5em, itemsep=2pt]
    \item At $\order(\varepsilon^0)$, $\lim_{\varepsilon\to 0^+}\psi(x,t+\varepsilon) = \psi(x,t)$ requires
    \begin{equation}
        A(\varepsilon) = I_0 = \sqrt{\frac{2\pi\ii\hbar\varepsilon}{m}}.
        \label{eq:A-fixed}
    \end{equation}
    \item At $\order(\varepsilon^1)$, subtracting $\psi(x,t)$, dividing by $\varepsilon$, taking $\varepsilon\to 0^+$, and multiplying by $\ii\hbar$ yields the Schr\"odinger equation:
    \begin{equation}
        \ii\hbar\frac{\partial\psi}{\partial t}
        = -\frac{\hbar^2}{2m}\frac{\partial^2\psi}{\partial x^2} + V(x)\psi.
    \end{equation}
\end{itemize}

\subsection{Time Slicing and the Path Integral}

For a finite time interval $T = t_f - t_i > 0$, divide $[t_i,t_f]$ into $N$ steps of length $\varepsilon = T/N$, with $t_k = t_i + k\varepsilon$, $x_0 = x_i$, and $x_N = x_f$. Iterating~\eqref{eq:feynman-ansatz} $N$ times gives $\psi(x_f,t_f) = \int K(x_f,t_f;\,x_i,t_i)\psi(x_i,t_i)\,\dd x_i$, where
\begin{equation}
    K(x_f,t_f;\,x_i,t_i)
    = \lim_{N\to\infty}\left(\frac{m}{2\pi\ii\hbar\varepsilon}\right)^{N/2}
      \int_{-\infty}^{\infty}\dd x_1 \cdots \int_{-\infty}^{\infty}\dd x_{N-1}\,
      \exp\!\left\{\frac{\ii}{\hbar}\sum_{k=0}^{N-1}\varepsilon\,
      \Lag\!\left(\bar x_k,\frac{x_{k+1}-x_k}{\varepsilon}\right)\right\},
    \label{eq:discretized-pi}
\end{equation}
with $\bar x_k := \frac{1}{2}(x_k+x_{k+1})$.

\begin{figure}[htbp]
    \centering
    \includegraphics[width=0.94\textwidth]{assets/fig3_time_slicing.pdf}
    \caption{Time-slicing discretization: each choice of $(x_1,\dots,x_{N-1})$ defines a polygonal path from $(t_i,x_i)$ to $(t_f,x_f)$, and the single-step phases add up to the Riemann sum of $S[x]/\hbar$.}
    \label{fig:time-slicing}
\end{figure}

As $N\to\infty$, the sum in the exponent approaches the action $S[x(t)] = \int_{t_i}^{t_f}\Lag(x,\dot x)\,\dd t$ along the path (Figure~\ref{fig:time-slicing}). Writing $\Dpath x(t)$ for the formal limit of the normalized integration measure gives the path-integral expression:

\begin{definitionbox}[Feynman Path Integral]
\begin{equation}
    K(x_f,t_f;\,x_i,t_i)
    = \int_{x(t_i)=x_i}^{x(t_f)=x_f} \Dpath x(t)\,\exp\!\left(\frac{\ii}{\hbar}S[x(t)]\right).
    \label{eq:feynman-pi}
\end{equation}
\end{definitionbox}

% ============================================================================
\section[Modern Operator Proof]{Modern Operator Proof: From $\hat U(t_f,t_i)$ to the Path Integral}
\label{sec:modern}
% ============================================================================

We now prove~\eqref{eq:discretized-pi}--\eqref{eq:feynman-pi} deductively from standard operator quantum mechanics without guessing the short-time kernel (Figure~\ref{fig:two-routes}).

\begin{figure}[htbp]
    \centering
    \includegraphics[width=0.96\textwidth]{assets/fig4_two_routes.pdf}
    \caption{Comparison of the two derivations: Feynman's historical route from $\Lag$ to the Schr\"odinger equation (top) and the modern operator derivation from $\hat H$ to the path integral (bottom).}
    \label{fig:two-routes}
\end{figure}

\subsection{Propagator as the Position Kernel of \texorpdfstring{$\hat U(t_f,t_i)$}{U(t\_f, t\_i)}}

Consider $\hat H = \hat T + \hat V = \frac{\hat p^2}{2m} + V(\hat x)$, with position and momentum eigenstates satisfying
\begin{equation}
    \braket{x'|x} = \delta(x'-x), \quad
    \braket{p'|p} = \delta(p'-p), \quad
    \int_{-\infty}^{\infty}\ket{x}\bra{x}\,\dd x = \hat I, \quad
    \int_{-\infty}^{\infty}\ket{p}\bra{p}\,\dd p = \hat I,
\end{equation}
and transformation kernel
\begin{equation}
    \braket{x|p} = \frac{1}{\sqrt{2\pi\hbar}}\,\ee^{\frac{\ii}{\hbar}px}, \qquad
    \braket{p|x} = \frac{1}{\sqrt{2\pi\hbar}}\,\ee^{-\frac{\ii}{\hbar}px}.
    \label{eq:xp-overlap}
\end{equation}

\begin{definitionbox}[Propagator]
For $t_f > t_i$ and $\hat U(t_f,t_i) = \ee^{-\ii\hat H(t_f-t_i)/\hbar}$, the propagator is the position-space kernel of $\hat U$:
\begin{equation}
    K(x_f,t_f;\,x_i,t_i) := \braket{x_f|\hat U(t_f,t_i)|x_i}
    = \braket{x_f|\ee^{-\frac{\ii}{\hbar}\hat H(t_f-t_i)}|x_i}.
\end{equation}
\end{definitionbox}

Inserting $\int\ket{x_i}\bra{x_i}\,\dd x_i = \hat I$ into $\braket{x_f|\psi(t_f)} = \braket{x_f|\hat U(t_f,t_i)|\psi(t_i)}$ gives
\begin{equation}
    \psi(x_f,t_f) = \int_{-\infty}^{\infty} K(x_f,t_f;\,x_i,t_i)\,\psi(x_i,t_i)\,\dd x_i.
\end{equation}
Three properties follow directly from $\hat U$:
\begin{enumerate}[leftmargin=1.5em, itemsep=2pt]
    \item \textbf{Initial limit:} $\lim_{t_f\to t_i^+} K(x_f,t_f;\,x_i,t_i) = \braket{x_f|\hat I|x_i} = \delta(x_f-x_i)$.
    \item \textbf{Differential equation:} $\bigl(\ii\hbar\partial_{t_f} - \hat H_{x_f}\bigr)K(x_f,t_f;\,x_i,t_i) = 0$ for $t_f > t_i$.
    \item \textbf{Composition:} From $\hat U(t_f,t_i) = \hat U(t_f,t_1)\hat U(t_1,t_i)$ and $\int\ket{x_1}\bra{x_1}\,\dd x_1 = \hat I$,
    \begin{equation}
        K(x_f,t_f;\,x_i,t_i) = \int_{-\infty}^{\infty} K(x_f,t_f;\,x_1,t_1)\,K(x_1,t_1;\,x_i,t_i)\,\dd x_1.
        \label{eq:composition-law}
    \end{equation}
\end{enumerate}

\subsection{Time Slicing and the Lie--Trotter Formula}

With $T = t_f - t_i$ and $\varepsilon = T/N$, writing $\hat U(T) = [\hat U(\varepsilon)]^N$ and inserting $(N-1)$ position resolutions gives
\begin{equation}
    K(x_f,t_f;\,x_i,t_i) = \int_{-\infty}^{\infty}\dd x_1\cdots\int_{-\infty}^{\infty}\dd x_{N-1}\,
    \prod_{k=0}^{N-1}\braket{x_{k+1}|\ee^{-\frac{\ii\varepsilon}{\hbar}(\hat T+\hat V)}|x_k}.
\end{equation}
Since $[\hat T,\hat V]\neq 0$, $\ee^{-\frac{\ii\varepsilon}{\hbar}(\hat T+\hat V)}$ does not factor exactly, but the Baker--Campbell--Hausdorff expansion gives
\begin{align}
    \ee^{-\frac{\ii\varepsilon}{\hbar}(\hat T+\hat V)}
    &= \ee^{-\frac{\ii\varepsilon}{\hbar}\hat T}\,\ee^{-\frac{\ii\varepsilon}{\hbar}\hat V} + \order(\varepsilon^2),
    \label{eq:trotter-1} \\
    \ee^{-\frac{\ii\varepsilon}{\hbar}(\hat T+\hat V)}
    &= \ee^{-\frac{\ii\varepsilon}{2\hbar}\hat V}\,\ee^{-\frac{\ii\varepsilon}{\hbar}\hat T}\,\ee^{-\frac{\ii\varepsilon}{2\hbar}\hat V} + \order(\varepsilon^3).
    \label{eq:trotter-2}
\end{align}

\begin{propositionbox}[Global Error Bound (Lie--Trotter)]
Let $\hat A_\varepsilon := \ee^{-\frac{\ii\varepsilon}{\hbar}(\hat T+\hat V)}$ and let $\hat B_\varepsilon$ denote the product in~\eqref{eq:trotter-1} ($r=2$) or~\eqref{eq:trotter-2} ($r=3$). Since $\|\hat A_\varepsilon\|=\|\hat B_\varepsilon\|=1$, the telescoping identity
\begin{equation}
    \hat A_\varepsilon^N - \hat B_\varepsilon^N
    = \sum_{j=0}^{N-1} \hat A_\varepsilon^{N-1-j}\,(\hat A_\varepsilon - \hat B_\varepsilon)\,\hat B_\varepsilon^{j}
\end{equation}
implies that on smooth states $\ket{\psi}$,
\begin{equation}
    \bigl\|(\hat A_\varepsilon^N - \hat B_\varepsilon^N)\ket{\psi}\bigr\|
    \le \sum_{j=0}^{N-1}\bigl\|(\hat A_\varepsilon - \hat B_\varepsilon)\hat B_\varepsilon^j\ket{\psi}\bigr\|
    = N\cdot\order(\varepsilon^r) = T\cdot\order(\varepsilon^{r-1}) \;\xrightarrow{\;N\to\infty\;}\; 0.
\end{equation}
More generally, if $\hat T+\hat V$ is essentially self-adjoint, $\hat B_{T/N}^N \to \ee^{-\ii(\hat T+\hat V)T/\hbar}$ strongly as $N\to\infty$.
\end{propositionbox}

\subsection{Momentum Completeness and the Legendre Transform}

Using the symmetric splitting~\eqref{eq:trotter-2}, the potential operators act on $\bra{x_{k+1}}$ and $\ket{x_k}$ as numbers:
\begin{equation}
    \braket{x_{k+1}|\hat B_\varepsilon|x_k}
    = \ee^{-\frac{\ii\varepsilon}{2\hbar}\bigl[V(x_{k+1})+V(x_k)\bigr]}
      \braket{x_{k+1}|\ee^{-\frac{\ii\varepsilon}{\hbar}\frac{\hat p^2}{2m}}|x_k}.
    \label{eq:pot-factored}
\end{equation}

Why insert a momentum completeness relation $\int\ket{p_k}\bra{p_k}\,\dd p_k = \hat I$ into the remaining factor?
\begin{enumerate}[leftmargin=1.5em, itemsep=2pt]
    \item \textbf{Diagonalizing $\hat T$:} In the position basis, $\hat p = -\ii\hbar\partial_x$ is a differential operator. Switching to $\ket{p_k}$ turns the operator $\ee^{-\ii\varepsilon\hat p^2/(2m\hbar)}$ into the number $\ee^{-\ii\varepsilon p_k^2/(2m\hbar)}$.
    \item \textbf{Legendre transform:} Integrating out $p_k$ replaces $p_k\dot x_k - H(p_k,x_k)$ in the exponent with $\Lag(x_k,\dot x_k)$.
\end{enumerate}

Inserting $\int\ket{p_k}\bra{p_k}\,\dd p_k = \hat I$ and using~\eqref{eq:xp-overlap} gives
\begin{equation}
\begin{aligned}
    \braket{x_{k+1}|\ee^{-\frac{\ii\varepsilon}{\hbar}\frac{\hat p^2}{2m}}|x_k}
    &= \int_{-\infty}^{\infty}\braket{x_{k+1}|p_k}\,\ee^{-\frac{\ii\varepsilon p_k^2}{2m\hbar}}\,\braket{p_k|x_k}\,\dd p_k \\
    &= \int_{-\infty}^{\infty}\frac{\dd p_k}{2\pi\hbar}\,
      \exp\!\left[\frac{\ii}{\hbar}p_k(x_{k+1}-x_k) - \frac{\ii\varepsilon}{\hbar}\frac{p_k^2}{2m}\right].
\end{aligned}
\label{eq:mom-inserted}
\end{equation}
Combining~\eqref{eq:pot-factored} and~\eqref{eq:mom-inserted} across all $N$ slices (with $H(p_k,\bar x_k) := \frac{p_k^2}{2m} + \frac{1}{2}[V(x_k)+V(x_{k+1})]$) before integrating over $p_k$ gives the Hamiltonian form of the path integral:

\begin{theorembox}[Phase-Space Path Integral]
\begin{equation}
\begin{aligned}
    K(x_f,t_f;\,x_i,t_i)
    = \lim_{N\to\infty}\int &\left(\prod_{k=1}^{N-1}\dd x_k\right)\left(\prod_{k=0}^{N-1}\frac{\dd p_k}{2\pi\hbar}\right) \\
      &\times \exp\!\left\{\frac{\ii}{\hbar}\sum_{k=0}^{N-1}\Bigl[p_k(x_{k+1}-x_k) - \varepsilon\,H(p_k,\bar x_k)\Bigr]\right\}.
\end{aligned}
\label{eq:phase-space-pi}
\end{equation}
\end{theorembox}

Now complete the square in $p_k$ inside~\eqref{eq:mom-inserted}, writing $\Delta x_k := x_{k+1}-x_k$:
\begin{equation}
    -\frac{\ii\varepsilon}{2m\hbar}p_k^2 + \frac{\ii}{\hbar}p_k\Delta x_k
    = -\frac{\ii\varepsilon}{2m\hbar}\!\left(p_k - m\frac{\Delta x_k}{\varepsilon}\right)^2
      + \frac{\ii\varepsilon}{\hbar}\cdot\frac{m}{2}\!\left(\frac{\Delta x_k}{\varepsilon}\right)^2.
\end{equation}
Evaluating the Gaussian integral over $u = p_k - m\Delta x_k/\varepsilon$ via~\eqref{eq:gaussian-moments} ($\alpha = \frac{\ii\varepsilon}{2m\hbar}+0^+$) gives
\begin{equation}
    \int_{-\infty}^{\infty}\frac{\dd u}{2\pi\hbar}\,\ee^{-\frac{\ii\varepsilon}{2m\hbar}u^2}
    = \frac{1}{2\pi\hbar}\sqrt{\frac{2\pi m\hbar}{\ii\varepsilon}}
    = \sqrt{\frac{m}{2\pi\ii\hbar\varepsilon}}.
\end{equation}
Restoring the potential factor yields the short-time propagator:
\begin{equation}
    K(x_{k+1},t_{k+1};\,x_k,t_k)
    = \sqrt{\frac{m}{2\pi\ii\hbar\varepsilon}}
      \exp\!\left\{\frac{\ii\varepsilon}{\hbar}\left[
          \frac{m}{2}\!\left(\frac{x_{k+1}-x_k}{\varepsilon}\right)^2
          - \frac{V(x_k)+V(x_{k+1})}{2}
      \right]\right\} + \order(\varepsilon^3).
    \label{eq:infinitesimal-propagator}
\end{equation}
At the stationary momentum $p_k^* = m\Delta x_k/\varepsilon = m\dot x_k$, the exponent $p_k^*\dot x_k - H(p_k^*,\bar x_k)$ becomes $\Lag(\bar x_k,\dot x_k)$, and the Gaussian integral supplies the prefactor $\sqrt{m/(2\pi\ii\hbar\varepsilon)}$ that had to be matched manually in~\eqref{eq:A-fixed}. Taking the product over $k=0,\dots,N-1$ and $N\to\infty$ recovers~\eqref{eq:discretized-pi}--\eqref{eq:feynman-pi}. \hfill$\square$

\subsection{Path Roughness and Midpoint Ordering}

\begin{remarkbox}[Nowhere-Differentiable Paths and Magnetic Vector Potentials]
\begin{enumerate}[leftmargin=1.5em, itemsep=3pt]
    \item \textbf{Roughness:} Dominant steps satisfy $|\Delta x_k|\sim\sqrt{\hbar\varepsilon/m}$, so typical paths are continuous ($\Delta x_k\to 0$) but have divergent velocity $|\Delta x_k/\varepsilon|\sim\varepsilon^{-1/2}\to\infty$, analogous to Brownian paths under $t = -\ii\tau$.
    \item \textbf{Midpoint rule:} For a scalar potential $V(x)$, shifting the evaluation point changes the exponent by $\order(\varepsilon\Delta x_k) = \order(\varepsilon^{3/2})$, which vanishes after $N=T/\varepsilon$ steps. For a magnetic vector potential term $q A(x)\dot x$, the discrete step is $q A(\tilde x_k)\Delta x_k$; shifting $\tilde x_k$ from $x_k$ to $\bar x_k = \frac{1}{2}(x_k+x_{k+1})$ changes the exponent by $\frac{\ii q}{2\hbar}A'(x_k)(\Delta x_k)^2 \sim \order(\varepsilon)$, leaving a finite $\order(1)$ difference as $N\to\infty$. Midpoint evaluation matches the Hermitian Weyl ordering $\frac{1}{2}[\hat p A(\hat x)+A(\hat x)\hat p]$ and preserves gauge invariance.
\end{enumerate}
\end{remarkbox}

% ============================================================================
\section[Worked Example: The Free Particle]{Worked Example: The Free Particle}
\label{sec:free}
% ============================================================================

\subsection{\texorpdfstring{$N$}{N}-Slice Gaussian Integration}

For a free particle ($V=0$, $T=t_f-t_i>0$), equation~\eqref{eq:discretized-pi} reads
\begin{equation}
    K_0(x_f,T;\,x_i,0) = \lim_{N\to\infty} \left(\frac{m}{2\pi\ii\hbar\varepsilon}\right)^{N/2}
    \int_{-\infty}^{\infty}\dd x_1\cdots\int_{-\infty}^{\infty}\dd x_{N-1}\,
    \exp\!\left[\frac{\ii m}{2\hbar\varepsilon}\sum_{k=0}^{N-1}(x_{k+1}-x_k)^2\right].
    \label{eq:free-N-slice}
\end{equation}

\begin{propositionbox}[Two-Step Gaussian Convolution]
For any $\varepsilon_1, \varepsilon_2 > 0$,
\begin{equation}
\begin{aligned}
    &\sqrt{\frac{m}{2\pi\ii\hbar\varepsilon_1}}\sqrt{\frac{m}{2\pi\ii\hbar\varepsilon_2}}
    \int_{-\infty}^{\infty}
    \exp\!\left[\frac{\ii m}{2\hbar}\!\left(\frac{(x_2-x_1)^2}{\varepsilon_2}+\frac{(x_1-x_0)^2}{\varepsilon_1}\right)\right]\dd x_1 \\[3pt]
    &\qquad\qquad = \sqrt{\frac{m}{2\pi\ii\hbar(\varepsilon_1+\varepsilon_2)}}
      \exp\!\left[\frac{\ii m(x_2-x_0)^2}{2\hbar(\varepsilon_1+\varepsilon_2)}\right].
\end{aligned}
\label{eq:gauss-conv-lemma}
\end{equation}
\end{propositionbox}

\textit{Proof.}\;
Set $a := \frac{1}{\varepsilon_1} + \frac{1}{\varepsilon_2} = \frac{\varepsilon_1+\varepsilon_2}{\varepsilon_1\varepsilon_2}$ and complete the square in $x_1$:
\begin{align}
    \frac{(x_2-x_1)^2}{\varepsilon_2} + \frac{(x_1-x_0)^2}{\varepsilon_1}
    &= a\!\left(x_1 - \frac{\varepsilon_1 x_2 + \varepsilon_2 x_0}{\varepsilon_1+\varepsilon_2}\right)^2
       + \frac{x_2^2}{\varepsilon_2} + \frac{x_0^2}{\varepsilon_1}
       - \frac{(\varepsilon_1 x_2 + \varepsilon_2 x_0)^2}{\varepsilon_1\varepsilon_2(\varepsilon_1+\varepsilon_2)} \nonumber\\[3pt]
    &= a\!\left(x_1 - \frac{\varepsilon_1 x_2 + \varepsilon_2 x_0}{\varepsilon_1+\varepsilon_2}\right)^2
       + \frac{(x_2 - x_0)^2}{\varepsilon_1 + \varepsilon_2}.
\end{align}
Integrating over $u = x_1 - \frac{\varepsilon_1 x_2 + \varepsilon_2 x_0}{\varepsilon_1+\varepsilon_2}$ via~\eqref{eq:gaussian-moments} gives $\sqrt{\frac{2\pi\ii\hbar}{ma}} = \sqrt{\frac{2\pi\ii\hbar\,\varepsilon_1\varepsilon_2}{m(\varepsilon_1+\varepsilon_2)}}$, which combines with the prefactors to give~\eqref{eq:gauss-conv-lemma}. \hfill$\square$

Applying~\eqref{eq:gauss-conv-lemma} to $x_1, x_2, \dots, x_{N-1}$ in turn replaces $(\varepsilon,\dots,\varepsilon)$ by $N\varepsilon = T$ at every finite $N$:

\begin{theorembox}[Free-Particle Propagator]
\begin{equation}
    K_0(x_f,T;\,x_i,0)
    = \sqrt{\frac{m}{2\pi\ii\hbar T}}\exp\!\left[\frac{\ii m(x_f - x_i)^2}{2\hbar T}\right]
    = \sqrt{\frac{m}{2\pi\ii\hbar T}}\exp\!\left(\frac{\ii}{\hbar}\Scl(x_f,T;\,x_i,0)\right),
    \label{eq:free-propagator}
\end{equation}
where $\Scl = \frac{m(x_f-x_i)^2}{2T}$ is the classical action of $\xcl(t) = x_i + \frac{x_f-x_i}{T}t$.
\end{theorembox}

\subsection{Consistency Checks}

\begin{enumerate}[leftmargin=1.5em, itemsep=3pt]
    \item \textbf{Momentum basis:} Inserting $\int\ket{p}\bra{p}\,\dd p = \hat I$ once on the full interval $T$ gives
    \[
        \braket{x_f|\ee^{-\frac{\ii T}{\hbar}\frac{\hat p^2}{2m}}|x_i}
        = \int_{-\infty}^{\infty}\frac{\dd p}{2\pi\hbar}\,\ee^{\frac{\ii}{\hbar}p(x_f-x_i)-\frac{\ii p^2 T}{2m\hbar}}
        = \sqrt{\frac{m}{2\pi\ii\hbar T}}\exp\!\left[\frac{\ii m(x_f-x_i)^2}{2\hbar T}\right].
    \]
    \item \textbf{Schr\"odinger equation:} Direct differentiation gives $\ii\hbar\partial_T K_0 = -\frac{\hbar^2}{2m}\partial_{x_f}^2 K_0 = \bigl[-\frac{\ii\hbar}{2T} + \frac{m(x_f-x_i)^2}{2T^2}\bigr]K_0$.
    \item \textbf{Initial condition:} Setting $u = (x_i-x_f)/\sqrt{\hbar T/m}$ gives
    \[
        \int_{-\infty}^{\infty} K_0(x_f,T;\,x_i,0)f(x_i)\,\dd x_i
        = \frac{1}{\sqrt{2\pi\ii}}\int_{-\infty}^{\infty} \ee^{\ii u^2/2}f\bigl(x_f + u\sqrt{\hbar T/m}\bigr)\,\dd u
        \;\xrightarrow{\;T\to 0^+\;}\; f(x_f),
    \]
    so $\lim_{T\to 0^+}K_0(x_f,T;\,x_i,0) = \delta(x_f-x_i)$.
\end{enumerate}

\subsection{Spreading of a Gaussian Wave Packet}
\label{sec:free-packet}

Take the initial Gaussian state with position width $\sigma_0$ and mean momentum $p_0 = m v_0$:
\begin{equation}
    \psi(y,0) = \frac{1}{(2\pi\sigma_0^2)^{1/4}}
    \exp\!\left(-\frac{y^2}{4\sigma_0^2} + \frac{\ii p_0 y}{\hbar}\right).
\end{equation}
By~\eqref{eq:free-propagator}, $\psi(x,T) = \int_{-\infty}^{\infty}K_0(x,T;\,y,0)\psi(y,0)\,\dd y$ reduces to the Gaussian integral $\int \ee^{-\gamma y^2 + \beta y}\,\dd y = \sqrt{\pi/\gamma}\,\ee^{\beta^2/(4\gamma)}$, where
\begin{equation}
    \gamma = \frac{1}{4\sigma_0^2} - \frac{\ii m}{2\hbar T} = \frac{m}{2\ii\hbar T}\!\left(1 + \frac{\ii T}{\tau_0}\right), \qquad
    \beta = \frac{\ii}{\hbar}\!\left(p_0 - \frac{m x}{T}\right), \qquad
    \tau_0 := \frac{2m\sigma_0^2}{\hbar}.
\end{equation}
Evaluating the integral and taking the modulus squared gives
\begin{equation}
    |\psi(x,T)|^2 = \frac{1}{\sqrt{2\pi}\,\sigma(T)}
    \exp\!\left[-\frac{(x - v_0 T)^2}{2\sigma(T)^2}\right],
    \qquad
    \sigma(T) = \sigma_0\sqrt{1 + \left(\frac{T}{\tau_0}\right)^2}
              = \sqrt{\sigma_0^2 + \left(\frac{\hbar T}{2m\sigma_0}\right)^2}.
    \label{eq:wavepacket-spreading}
\end{equation}
An initial spread $\Delta x_0 = \sigma_0$ implies $\Delta p_0 = \hbar/(2\sigma_0)$, or $\Delta v_0 = \hbar/(2m\sigma_0)$. Thus the probability band follows $x_{\mathrm{cl}}(T) = v_0 T$ while widening as $\sigma(T) \approx (\Delta v_0)\,T$ for $T \gg \tau_0$ (Figure~\ref{fig:free-wavepacket}).

\begin{figure}[htbp]
    \centering
    \includegraphics[width=0.96\textwidth]{assets/fig5_free_wavepacket.pdf}
    \caption{Gaussian wave-packet evolution from~\eqref{eq:wavepacket-spreading} ($\sigma_0=1$, $\tau_0=1$, $v_0=2\sigma_0/\tau_0$). \textbf{(a)} Mean trajectory $\langle\hat x\rangle=v_0 t$ and $\pm\sigma(t)$ envelope. \textbf{(b)} Density profiles $|\psi(x,t)|^2$ at three times.}
    \label{fig:free-wavepacket}
\end{figure}

% ============================================================================
\section[Stationary Phase \& Corridor Width]{Stationary Phase and the Width of the Classical Corridor}
\label{sec:semiclassical}
% ============================================================================

\subsection{Expansion Around the Classical Path}

Write any path from $(t_i,x_i)$ to $(t_f,x_f)$ as $x(t) = \xcl(t) + y(t)$ with $y(t_i)=y(t_f)=0$. Expanding $S[\xcl+y]$ in $y(t)$ and integrating the linear term by parts gives
\begin{equation}
    S[\xcl + y]
    = S[\xcl]
    + \underbrace{\int_{t_i}^{t_f} \left[\frac{\partial\Lag}{\partial x} - \frac{\dd}{\dd t}\frac{\partial\Lag}{\partial\dot x}\right]_{\xcl} y(t)\,\dd t}_{=\,0\;\text{(Euler--Lagrange)}}
    \;+\; \frac{1}{2}\delta^2 S[y] + \order(y^3),
\end{equation}
where $\frac{1}{2}\delta^2 S[y] = \int_{t_i}^{t_f}\bigl[\frac{1}{2}m\dot y^2 - \frac{1}{2}V''(\xcl(t))y^2\bigr]\dd t$. Shifting integration variables from $x(t)$ to $y(t)$ factors out $\ee^{\ii S[\xcl]/\hbar}$:
\begin{equation}
    K(x_f,t_f;\,x_i,t_i)
    = \ee^{\frac{\ii}{\hbar}S[\xcl]}
      \int_{y(t_i)=0}^{y(t_f)=0} \Dpath y(t)\,
      \exp\!\left(\frac{\ii}{2\hbar}\delta^2 S[y] + \frac{\ii}{\hbar}\order(y^3)\right).
    \label{eq:fluctuation-pi}
\end{equation}
When $V(x)$ is at most quadratic, $\order(y^3)=0$ and the $y$-integral depends only on $t_f, t_i$, so $K$ has the exact form $F(T)\,\ee^{\ii\Scl/\hbar}$ as seen in~\eqref{eq:free-propagator}.

For general potentials with $S \gg \hbar$, adding the unit phasors $\ee^{\ii S[x]/\hbar}$ in the complex plane produces two behaviors (Figure~\ref{fig:stationary-phase}):
\begin{itemize}[leftmargin=1.5em, itemsep=2pt]
    \item Away from $\xcl$ ($\delta S \neq 0$), small path changes rotate $S/\hbar$ rapidly, winding the phasor sum into self-cancelling spirals.
    \item Near $\xcl$ ($\delta S = 0$), $S[x]$ is stationary, so paths satisfying $|S[x]-S[\xcl]|\lesssim\hbar$ have aligned phases and produce the dominant contribution to $K$. Thus $\delta S = 0$ arises in classical mechanics as the stationary-phase condition of quantum interference.
\end{itemize}

\begin{figure}[htbp]
    \centering
    \includegraphics[width=0.96\textwidth]{assets/fig6_stationary_phase.pdf}
    \caption{\textbf{(a)} Action $S[\xcl+y]/\hbar$ near $y=0$ ($\delta S=0$), showing the stationary band $\Delta S\lesssim\hbar$. \textbf{(b)} Complex phasor sum (Cornu spiral): far paths curl up and cancel, while paths near $\xcl$ add coherently.}
    \label{fig:stationary-phase}
\end{figure}

\subsection{Estimating the Corridor Width}

For a free particle traveling distance $L = x_f - x_i$ in time $T$ ($v = L/T$), take a piecewise-linear detour $y(t)$ reaching $\Delta y$ at $t = T/2$ and vanishing at $t=0, T$:
\begin{equation}
    y(t) =
    \begin{cases}
        \dfrac{2\Delta y}{T}\,t, & 0 \le t \le \dfrac{T}{2}, \\[6pt]
        \dfrac{2\Delta y}{T}(T-t), & \dfrac{T}{2} \le t \le T.
    \end{cases}
\end{equation}
The excess action is $\Delta S = \int_0^T \frac{1}{2}m\dot y^2\,\dd t = \frac{2m(\Delta y)^2}{T}$. Setting $\Delta S \sim \hbar$ gives the corridor half-width
\begin{equation}
    \Delta y_{\mathrm{tube}} \sim \sqrt{\frac{\hbar T}{2m}}
    = \sqrt{\frac{\lambda_{\mathrm{dB}}\,L}{4\pi}},
    \qquad
    \lambda_{\mathrm{dB}} := \frac{h}{mv},
    \label{eq:tube-width}
\end{equation}
which has the same $\sqrt{\lambda L}$ scaling as the first Fresnel zone in wave optics. Table~\ref{tab:tube-scales} compares~\eqref{eq:tube-width} for a macroscopic body and a low-energy electron.

\begin{table}[!htbp]
    \centering
    \small
    \setlength{\tabcolsep}{6pt}
    \renewcommand{\arraystretch}{1.25}
    \begin{tabular}{@{}l c c c >{\raggedright\arraybackslash}p{4.4cm}@{}}
        \toprule
        \textbf{System} & \textbf{Mass $m$} & \textbf{Flight $(L, T)$} & \textbf{Width $\Delta y_{\mathrm{tube}}$} & \textbf{Regime} \\
        \midrule
        Marble & $1.0\,\text{g}$ & $1\,\text{m},\; 1\,\text{s}$ & $2.3 \times 10^{-16}\,\text{m}$ & Sub-nuclear; sharp classical curve \\
        $1\,\text{eV}$ electron & $9.11 \times 10^{-31}\,\text{kg}$ & $1\,\text{cm},\; 17\,\text{ns}$ & $9.9 \times 10^{-7}\,\text{m}$ & $\sim 1\,\mu\text{m}$; wide diffraction regime \\
        \bottomrule
    \end{tabular}
    \caption{Characteristic corridor width $\Delta y_{\mathrm{tube}} = \sqrt{\hbar T/(2m)}$ in two regimes.}
    \label{tab:tube-scales}
\end{table}

\end{document}
